\chapter{Agenda et planning}
\label{ch:agenda}

Le module \textbf{Agenda} sert à saisir, modifier et visualiser les \textbf{sessions} d'intervention. En v2, la préparation de facturation n'y figure plus~: elle est sur la fiche école (chapitre~\ref{ch:facturation}).

\section{Accès}

\menu{Agenda} dans la barre latérale → \texttt{/planning}. Deux onglets~: \menu{Planning} et \menu{Calendrier}. Bascule \textbf{mois} / \textbf{semaine} (mémorisée en session).

\screenshot[keepaspectratio]{05-planning.png}{Vue planning — sélection école et cours, calendrier mensuel}

\section{Sélection du contexte}

Avant de planifier, choisissez~:

\begin{enumerate}[label=\textbf{\arabic*.}]
  \item L'\textbf{école} (établissement client).
  \item Le \textbf{cours} concerné — le sélecteur du fil d'Ariane liste \textbf{tous} les cours de l'école, y compris ceux d'une autre année civile.
  \item Le \textbf{groupe} si plusieurs groupes sont liés au cours.
\end{enumerate}

Ces choix sont mémorisés en session pour accélérer les saisies successives.

\section{Vues mois et semaine}

\begin{itemize}[leftmargin=*]
  \item \textbf{Mois} (défaut)~: grille 7 colonnes sur le mois calendaire. Les flèches avancent d'un mois.
  \item \textbf{Semaine}~: lundi--dimanche, grille horaire 08{,}00--20{,}00. Les flèches avancent d'une semaine. Les événements hors plage sont masqués ou calés.
\end{itemize}

Les cartes KPI portent sur la \textbf{période visible} (semaine ou mois), montants en TTC (HT ×\,1{,}2). L'agenda liste les sessions par \textbf{date de début}, pas par l'année du cours~: un cours 2026 avec une séance en janvier 2027 apparaît bien en janvier 2027.

\section{Planifier une session}

\begin{workflowstep}[Depuis la grille]
  Mode \textbf{Édition} \emph{et} cours sélectionné. Sans cours, les numéros de jour restent inactifs. Cliquez le numéro du jour (ou l'en-tête de jour en vue semaine) pour créer à \textbf{08{,}00}, ou un créneau horaire (08{h}--19{h}) pour créer à cette heure pile.
\end{workflowstep}

\begin{workflowstep}[Formulaire]
  Complétez horaires, groupe, lieu. Validez avec \menu{Planifier}. Le bouton doit appartenir au même formulaire que les champs date/heure.
\end{workflowstep}

La durée facturable est calculée à partir des horaires réels. Un multiplicateur corrige les imports calendrier erronés (\texttt{billableMultiplier}).

\section{Dupliquer une session}

Depuis une case jour, un événement semaine ou la fiche session (mode Édition)~: recopier vers \textbf{demain}, \textbf{la semaine suivante} ou une \textbf{date personnalisée} (même heure, durée conservée).

\begin{itemize}[leftmargin=*]
  \item Bloqué si la session source a déjà un n° de facture.
  \item Refusé s'il existe déjà une session du même groupe sur le créneau cible.
  \item La date personnalisée s'ouvre dans une boîte de dialogue native (pas une modale Alpine).
\end{itemize}

\section{Navigation temporelle}

\begin{itemize}[leftmargin=*]
  \item Flèches précédent / suivant (mois ou semaine selon la vue).
  \item Clic sur une session existante pour l'éditer ou la dupliquer.
\end{itemize}

\section{Import calendrier (ICS)}

L'onglet \menu{Calendrier} (\texttt{/admin/calendars}) importe un fichier \texttt{.ics} \emph{ou} une URL ICS vers une \textbf{école}.

\begin{workflowstep}[Analyser puis mapper]
  Choisissez l'école, déposez le fichier (ou l'URL), \menu{Analyser}. Pour chaque libellé détecté dans le \textbf{SUMMARY}, associez un \textbf{cours et un groupe}, puis validez. Les sessions sont créées avec \texttt{billable\_rate = 1} (durée complète). Les chevauchements du même groupe sont ignorés.
\end{workflowstep}

\begin{itemize}[leftmargin=*]
  \item Cours \emph{et} groupe sont tous les deux requis pour créer une session, même si l'UI libelle le groupe «~optionnel~».
  \item Le \textbf{réimport} relit uniquement le mapping mémorisé (un morph cours \emph{ou} groupe, le cours l'emporte) et le champ SUMMARY~: après un mapping cours, le réimport ignore en pratique tous les événements faute de groupe.
  \item Supprimer une source d'import n'efface pas les sessions déjà créées (\texttt{calendar\_source\_id} n'est pas dans \texttt{Planning::\$fillable}).
\end{itemize}

Détail technique (wiki)~: \texttt{documentation/fr/import-calendrier.md}.

\section{Lien avec la facturation}

Chaque session planifiée alimente~:

\begin{itemize}[leftmargin=*]
  \item le montant \textbf{non facturé TTC} sur la liste des écoles~;
  \item la section \texttt{\#billing} de la fiche école pour le mois concerné~;
  \item l'histogramme \textbf{planifié} de la trésorerie.
\end{itemize}

La facturation elle-même se prépare depuis la fiche école, pas depuis l'agenda.
